\documentclass[11pt]{article}
\usepackage[T1]{fontenc}
\usepackage{lmodern}
\usepackage[margin=1in]{geometry}
\usepackage{setspace}
\usepackage{booktabs}
\usepackage{url}
\usepackage[round]{natbib}
\input{glyphtounicode}
\pdfgentounicode=1
\onehalfspacing
\hyphenpenalty=10000
\exhyphenpenalty=10000
\sloppy
\newcommand{\etap}{\ensuremath{\eta^2_p}}

\title{Looking Up, Asking Less: Status Cues and Help-Seeking at Work\\[1em]
\normalsize\textbf{SYNTHETIC TEST PAPER, NOT REAL RESEARCH}}
\author{Mara Vellacott\thanks{This is a synthetic test paper written to test a review tool. It is not real research, and none of its data exist. The authors, Eastmoor University and its Behavioral Lab, the Crestline Research Panel, and Norrland Freight are fictional. Any resemblance to real people or organizations is unintended.} \and Tobias Renquist\\[0.5em]
\small Eastmoor University}
\date{}

\begin{document}
\maketitle

\begin{abstract}
\noindent Employees often need help from colleagues who outrank them. We argue that a cue signaling a potential helper's high status reduces help-seeking from that helper. In an online experiment with 412 employed adults and a preregistered laboratory experiment with 240 participants, people sent fewer requests for help to a helper described as a senior executive than to one described as a peer. The effect was weaker when requests were anonymous, and evaluation apprehension carried it. A field test at a logistics firm, in which 30 of 60 teams were invited to an anonymous help platform, shows that using the platform raised supervisor-rated performance. We discuss what these results mean for how organizations display rank.

\medskip
\noindent\textit{Keywords:} help-seeking; status; evaluation apprehension; anonymity; experiments
\end{abstract}

\section{Introduction}

Much of the knowledge an employee needs sits in other people's heads. Asking for help is how that knowledge moves, and organizations invest in mentoring programs, expert directories, and internal question forums on the assumption that people will use them. Yet the person best placed to answer a question is often more senior than the person who has it. A new analyst whose model will not reconcile may know that the planning director built the original, and still decide to spend an afternoon guessing rather than send a message upstairs.

This paper asks whether a simple cue to a helper's status changes whether people ask that helper for help. Organizations make status visible all the time: titles in email signatures, org charts in collaboration software, seniority badges in internal chat tools. If such cues suppress requests, then some of the knowledge that sits with senior people never reaches the employees who need it, and the cost is paid quietly, by people who do not ask.

We make four contributions. First, we show in two experiments that attaching a high-status title to a potential helper reduces help-seeking from that helper, holding the helper's availability and expertise constant. Second, we show that the effect is weaker when requests are anonymous, which points to the requester's concern about how the helper will see them. Third, we show that evaluation apprehension, the worry of being judged as less capable, carries the effect. Fourth, we find that status affects the decision to ask but leaves the content of requests unchanged, which narrows where interventions should aim. A field test at a logistics firm then asks whether removing the evaluative sting of a request, by giving employees an anonymous channel, pays off in performance.

\section{Theory and Hypotheses}

\subsection{Status and the decision to ask}

A request for help discloses a gap. The requester reveals that there is something they do not know or cannot do, and the recipient of the request learns it. Whether that disclosure is costly depends on who receives it. A peer who learns of the gap has limited say over the requester's standing; a senior colleague may shape assignments, ratings, and promotions, and may form an impression that lasts. We therefore expect people to weigh the same request differently depending on the status of the person who would receive it.

Status cues are cheap to observe and hard to ignore. A title such as senior director or vice president tells a potential requester that the helper sits above them and may be consulted about them. Even when the helper is equally available and equally expert, the requester has more to lose by revealing a gap to that person. People can manage that risk by not asking.

The argument has a scope condition. It concerns help that the requester could obtain from more than one source, or could forgo at some cost to the work. When only one person can answer a question and the work cannot proceed without the answer, the requester has little room to choose, and the status of the helper should matter less for whether the request is made. Most everyday help at work is not of this kind. Employees can guess, search, ask a peer who knows less, or wait, and each of these options lets them avoid a senior colleague at some cost in quality or time. It is in this common situation that we expect status cues to shape the decision to ask.

\textit{Hypothesis 1 (H1). A cue signaling that a potential helper holds high status reduces help-seeking from that helper.}

\subsection{Anonymity}

If the cost of asking lies in what the helper learns about the requester, then removing the requester's identity should shrink the cost. An anonymous request still discloses a gap, but the gap cannot be attached to a person. The status of the recipient should then matter less, because the recipient can no longer form an impression of the particular requester.

\textit{Hypothesis 2 (H2). The negative effect of a high-status cue on help-seeking is weaker when requests are anonymous than when they are identified.}

\subsection{Evaluation apprehension}

The argument so far rests on a specific psychological state: the concern that one will be judged, and judged as less capable, by the person one asks. We call this evaluation apprehension. A high-status helper should raise it, because such a helper's judgment carries more weight, and higher apprehension should in turn lower the willingness to ask.

\textit{Hypothesis 3 (H3). Evaluation apprehension mediates the effect of a high-status cue on help-seeking.}

\subsection{An anonymous channel in the field}

If evaluation apprehension keeps people from asking the colleagues who could help them most, then a channel that lets employees ask without being identified should release some of the suppressed requests. When those requests reach people who can answer them, the employees who ask should do their work better.

\textit{Hypothesis 4 (H4). Using an anonymous help platform raises supervisor-rated performance.}

\section{Overview of Studies}

Study 1, an online experiment with employed adults, tests H1 and H2 in a 2 by 2 design that crosses the helper's status with the anonymity of requests. Study 2, a laboratory experiment registered before data collection, tests H1 with a different status cue and a behavioral measure of requests, and tests H3 by measuring evaluation apprehension. Study 3, a field test at a logistics firm, tests H4 by inviting randomly drawn teams to an anonymous help platform. Table~\ref{tab:flow} gives the flow of participants through each study. All three studies were approved by the Eastmoor University review board. In Studies 1 and 2 the helper's replies were prepared in advance, and every participant was told so in a debriefing at the end of the session.

\begin{table}[ht]
\centering
\caption{Participant flow}
\label{tab:flow}
\small
\begin{tabular}{llrrr}
\toprule
Study & Condition & Randomized & Excluded & Analyzed \\
\midrule
Study 1 & Senior director, identified & 103 & 5 & 98 \\
 & Senior director, anonymous & 103 & 5 & 98 \\
 & Peer, identified & 103 & 5 & 98 \\
 & Peer, anonymous & 103 & 5 & 98 \\
 & Total & 412 & 20 & 392 \\
\midrule
Study 2 & Vice president & 120 & 31 & 89 \\
 & Peer & 120 & 9 & 111 \\
 & Total & 240 & 40 & 200 \\
\midrule
Study 3 & Invited teams (30), employees & 596 & 0 & 596 \\
 & Teams not invited (30), employees & 588 & 0 & 588 \\
 & Total & 1,184 & 0 & 1,184 \\
\bottomrule
\end{tabular}

\smallskip
\parbox{0.9\linewidth}{\footnotesize\textit{Note.} Study 1 excluded participants who failed an attention item. Study 2 excluded participants who answered a comprehension question incorrectly. In Study 3, teams were randomized and employees analyzed.}
\end{table}

\section{Study 1}

\subsection{Method}

\paragraph{Participants.} We recruited employed adults in the United States from the Crestline Research Panel in March 2025 and paid each \$3.00 for a session of about 15 minutes. An a priori power analysis for a main effect of \textit{d} = 0.30, with 80 percent power and a two-sided alpha of .05, required 352 participants \citep{lakens2022}. We recruited 412, allowing for exclusions. The survey software assigned each participant to one of the four cells with equal allocation, 103 per cell; no researcher took part in assignment. Twenty participants failed an attention item placed after the manipulation, five per cell, leaving 392 participants, 98 per cell.\footnote{Study 1 was not preregistered. We set the attention rule before collecting data. Without the exclusion, the main effect of the status cue on all 412 randomized participants is {\itshape F}(1, 408) = 17.95, \textit{p} < .001, \etap{} = .042, and the status-by-request interaction is {\itshape F}(1, 408) = 3.30, \textit{p} = .070, \etap{} = .008.} A sensitivity analysis shows that 98 per cell gives 80 percent power to detect an interaction in which the status effect differs between request types by \textit{d} = 0.57, so the study is powered for the main effect in H1 and only for large interactions. All 412 participants completed the session. The panel removed duplicate devices before the study opened; we did not remove anyone for this reason.

\paragraph{Design and procedure.} The design was a 2 (status cue: senior director or peer) by 2 (request: anonymous or identified) between-subjects design. Participants were told that the study concerned how people solve estimation problems at work when a colleague is available. They read a short profile of a colleague, Dana Whitlock, who would be available during the task; the profile described her either as a senior director or as a planning analyst at the participant's own level (Appendix A). Participants in the identified condition learned that their first name would appear on any request they sent; those in the anonymous condition learned that Dana would see only the text of the request. After the profile, participants answered the attention item, then worked through 12 estimation problems. Each problem had a button labeled ``Ask Dana for a hint.'' A hint arrived about 20 seconds after a request.

\paragraph{Measures.} The outcome was the number of problems, out of 12, on which the participant asked Dana for a hint. After the task, participants rated how senior Dana was relative to them (1 = much more junior, 7 = much more senior) and how likely it was that Dana knew who had sent the requests (1 = certain she knew, 7 = certain she did not). No covariates were measured or used.

\subsection{Results}

We analyzed the number of hint requests in a two-way analysis of variance. Table~\ref{tab:s1} reports cell means and standard deviations.

\begin{table}[ht]
\centering
\caption{Study 1: Hint requests by condition}
\label{tab:s1}
\small
\begin{tabular}{lcccc}
\toprule
 & \multicolumn{2}{c}{Identified} & \multicolumn{2}{c}{Anonymous} \\
\cmidrule(lr){2-3}\cmidrule(lr){4-5}
Status cue & \textit{M} & \textit{SD} & \textit{M} & \textit{SD} \\
\midrule
Peer & 5.40 & 2.37 & 5.64 & 2.36 \\
Senior director & 3.92 & 2.44 & 5.04 & 2.43 \\
\midrule
\multicolumn{5}{l}{Analysis of variance} \\
\midrule
Source & {\itshape F}(1, 380) & \textit{p} & \etap & \\
\midrule
Status cue & 18.40 & .000 & .046 & \\
Request type & 7.84 & .005 & .020 & \\
Status cue by request type & 3.21 & .03 & .008 & \\
\bottomrule
\end{tabular}

\smallskip
\parbox{0.8\linewidth}{\footnotesize\textit{Note.} \textit{n} = 98 per cell. The outcome is the number of problems, out of 12, on which the participant asked for a hint. The software prints \textit{p} values below .0005 as .000.}
\end{table}

\paragraph{Manipulation checks.} Participants rated the senior director as more senior than the peer (\textit{M} = 6.12 versus 3.05), {\itshape F}(1, 380) = 412.36, \textit{p} < .001, \etap{} = .52. Participants in the anonymous condition were more confident that Dana did not know who had asked (\textit{M} = 5.88 versus 2.71), {\itshape F}(1, 380) = 226.50, \textit{p} < .001, \etap{} = .37.

\paragraph{Main effect of status.} Participants asked the senior director for fewer hints than the peer (\textit{M} = 4.48 versus 5.52), {\itshape F}(1, 380) = 18.40, \textit{p} = .000, \etap{} = .046, \textit{d} = 0.43. H1 is supported. Participants also asked for more hints when requests were anonymous (\textit{M} = 5.34 versus 4.66), {\itshape F}(1, 380) = 7.84, \textit{p} = .005, \etap{} = .020.

\paragraph{Anonymity as a moderator.} The status effect was larger when requests were identified (5.40 versus 3.92) than when they were anonymous (5.64 versus 5.04). The interaction was significant, {\itshape F}(1, 380) = 3.21, \textit{p} = .03, \etap{} = .008. H2 is supported.

\subsection{Discussion}

Study 1 shows that describing an available colleague as a senior director reduced how often people asked her for help, and that the reduction was smaller when requests were anonymous. Study 2 moves to the laboratory, uses a different status cue, adds a behavioral count of messages sent in a live chat, and measures evaluation apprehension to test H3.

\section{Study 2}

\subsection{Method}

\paragraph{Participants.} We recruited 240 students and staff from the participant pool of the Eastmoor University Behavioral Lab between February 10 and April 25, 2025. Each received \$10.00. The study was registered before data collection, which keeps its predictions separate from later analytic choices \citep{nosek2018}; Appendix B reproduces the registration.

\paragraph{Design and procedure.} Participants were randomly assigned to the vice-president or the peer condition, 120 per condition. Sessions ran on lab computers, and the experimenter, who greeted participants and started the software, did not know their condition. Participants first reported their age and completed a five-item prior help-seeking habit scale written for this study ($\alpha$ = .84; items in Appendix A). They then read a profile of Dana Whitlock, who would be available by chat during a 25-minute spreadsheet task. The profile described her either as the vice president of operations or as an operations analyst at the participant's level (Appendix A). Participants next completed a four-item help-seeking intention scale ($\alpha$ = .91), and every participant drafted the message they would send if they asked Dana for help. They then did the task, during which they could send Dana as many chat messages as they wished. After the task, participants answered a comprehension question about the task's rules, completed a four-item evaluation apprehension scale ($\alpha$ = .88), and rated Dana's seniority.

\paragraph{Measures.} The primary outcome, as registered, was the help-seeking intention scale (1 to 7). The behavioral outcome was the number of help requests sent to Dana during the task. The quality of help requested was the mean of two coders' ratings of each drafted message (1 to 5; intraclass correlation = .81); the coders did not know the participants' conditions. The two covariates, age and prior help-seeking habit, were measured before the manipulation.

\paragraph{Exclusions.} As preregistered, we excluded participants who answered the comprehension question incorrectly: 31 in the vice-president condition and 9 in the peer condition. The analyzed sample is 200 participants, 89 in the vice-president condition and 111 in the peer condition.

\paragraph{Analysis.} As registered, every outcome was analyzed with condition and the two covariates: an analysis of covariance for the intention scale and the quality rating, and a regression for the number of requests.

\subsection{Results}

\begin{table}[ht]
\centering
\caption{Study 2: Outcomes by condition, adjusted for age and prior help-seeking habit}
\label{tab:s2}
\small
\begin{tabular}{lccccc}
\toprule
Outcome & Peer & Vice president & Test & \textit{p} & \textit{d} \\
\midrule
Help-seeking intention (primary) & 4.63 & 4.28 & {\itshape F}(1, 196) = 5.02 & .026 & 0.32 \\
Help requests sent & 3.44 & 2.90 & \textit{t}(196) = 1.80 & .037 (one-tailed) & 0.26 \\
Quality of help requested & 3.12 & 3.23 & {\itshape F}(1, 196) = 0.96 & .33 & 0.14 \\
\bottomrule
\end{tabular}

\smallskip
\parbox{0.9\linewidth}{\footnotesize\textit{Note.} \textit{n} = 111 (peer) and 89 (vice president). Entries in the Peer and Vice president columns are adjusted means. \textit{d} is the standardized difference between conditions computed from the adjusted test statistic. The help-requests test is one-tailed, as registered; the other tests are two-tailed.}
\end{table}

\paragraph{Manipulation check.} Participants rated the vice president as more senior than the peer, {\itshape F}(1, 196) = 138.20, \textit{p} < .001, \etap{} = .41.

\paragraph{Help-seeking.} Table~\ref{tab:s2} reports the results. Participants in the vice-president condition reported lower help-seeking intention than those in the peer condition (adjusted \textit{M} = 4.28 versus 4.63), {\itshape F}(1, 196) = 5.02, \textit{p} = .026, \etap{} = .025, \textit{d} = 0.32, 95 percent CI 0.04 to 0.60. They also sent fewer help requests during the task (adjusted \textit{M} = 2.90 versus 3.44), \textit{t}(196) = 1.80, \textit{p} = .037, one-tailed, as preregistered, \textit{d} = 0.26. H1 is supported again.

\paragraph{Quality of help requested.} The messages drafted in the vice-president condition were rated slightly higher in quality than those in the peer condition, \textit{d} = 0.14, 95 percent CI $-$0.14 to 0.42, {\itshape F}(1, 196) = 0.96, \textit{p} = .33.

\paragraph{Mediation.} We estimated the indirect effect of condition on help-seeking intention through evaluation apprehension with 5,000 bootstrap resamples, coding the peer condition as 1 and including both covariates. The peer condition lowered evaluation apprehension (\textit{a} = $-$0.41, \textit{SE} = 0.13), and evaluation apprehension lowered help-seeking intention (\textit{b} = $-$0.22, \textit{SE} = 0.07). The indirect effect was 0.09, 95 percent bootstrap CI .02 to .19, and the direct effect was 0.26. This indirect effect establishes evaluation apprehension as the mechanism through which a high-status cue reduces help-seeking. H3 is supported.

\subsection{Discussion}

Study 2 replicates the effect of a status cue on help-seeking with a different title, in a different population, with both a registered intention measure and a count of messages actually sent. It also identifies the psychological state that carries the effect. Study 3 asks whether an intervention built on that account, an anonymous channel for asking, changes performance in a firm.

\section{Study 3}

\subsection{Setting and design}

Norrland Freight is a regional logistics firm whose dispatch and warehouse employees work in teams of 12 to 30 people, each team based at one depot and led by one supervisor. In September 2025 the firm introduced an internal platform on which employees could post work questions without their names and receive answers from colleagues across the firm, including senior staff. We designed the draw; the firm's HR analytics group carried it out. The group sorted the 60 team codes by a random number generated from a seed we supplied and invited the first 30 teams. Every employee on an invited team received an email invitation and a login; employees on other teams could reach the platform only through a link that circulated informally. The study was not preregistered.

The sample is the 1,184 employees on the 60 team rosters when the invitations went out, 596 on invited teams and 588 on the others. Rosters are fixed for the quarter, and the firm's quarter-end rating covered every one of them. The outcome is the supervisor's quarter-end performance rating, standardized to a mean of 0 and a standard deviation of 1. The treatment is use of the platform, recorded from its logs as posting at least one question in the twelve weeks after the invitations. Because the invitation was randomized by team, we report standard errors clustered by team in every model \citep{abadie2023, bruhn2009}.

We use the invitation as an instrument for platform use. We assume that the invitation affects performance ratings only through platform use. Because an invitation added a route to the platform and removed none, we also assume that no employee was less likely to use the platform because their team was invited.

\subsection{Results}

\begin{table}[ht]
\centering
\caption{Study 3: Invitation, platform use, and supervisor ratings}
\label{tab:s3}
\small
\begin{tabular}{lccc}
\toprule
 & (1) First stage & (2) Reduced form (ITT) & (3) 2SLS \\
Dependent variable & Used platform & Rating & Rating \\
\midrule
Invited & 0.20 & 0.09 & \\
 & & (0.04) & \\
Used platform & & & 0.45 \\
 & & & (0.20) \\
\midrule
Cragg-Donald {\itshape F} & 108.2 & & \\
Employees & 1,184 & 1,184 & 1,184 \\
Teams & 60 & 60 & 60 \\
\bottomrule
\end{tabular}

\smallskip
\parbox{0.85\linewidth}{\footnotesize\textit{Note.} Linear models with no controls. Standard errors clustered by team in parentheses. Rating is the supervisor's quarter-end performance rating in standard deviation units. In column 3, platform use is instrumented with the invitation.}
\end{table}

\paragraph{First stage.} In invited teams, 143 of 596 employees (24.0 percent) used the platform; in the other teams, 24 of 588 (4.1 percent) did. The invitation raised the probability of use by 0.20 (Table~\ref{tab:s3}, column 1). The Cragg-Donald {\itshape F} statistic is 108.2, far above the Stock-Yogo critical values, so the invitation is a strong instrument.

\paragraph{Effect on performance.} Column 2 reports the reduced form: the intention-to-treat effect of the invitation on supervisor ratings is 0.09 standard deviations (\textit{SE} = 0.04; 95 percent CI 0.01 to 0.17). Column 3 reports the two-stage least squares estimate: using the platform raised ratings by 0.45 standard deviations (\textit{SE} = 0.20; 95 percent CI 0.06 to 0.84), which equals the reduced form divided by the first stage (0.09 / 0.20). This is a local estimate. It applies to the compliers, the employees who used the platform because their team was invited and would not have used it otherwise, about one in five invited employees. It does not describe employees who would have found the platform without an invitation or who ignored it. H4 is supported.

\section{General Discussion}

Across two experiments, a cue signaling that an available colleague held high status reduced help-seeking from that colleague. In Study 1 the effect was weaker when requests were anonymous; in Study 2 evaluation apprehension carried it. Study 3 showed that giving employees an anonymous channel for their questions raised the performance ratings of those it drew in.

\subsection{Contributions}

The first contribution is the effect itself. Holding the helper's availability and expertise constant, a title was enough to reduce requests, in an online sample of employed adults and in a laboratory sample, with a registered intention measure and a count of messages sent.

The second contribution concerns the content of requests. One might expect people facing a senior helper to ask less often but to ask better, polishing their questions to protect their image. Study 2 found no such trade. The quality result (\textit{d} = 0.14, 95 percent CI $-$0.14 to 0.42) shows that status changes whether people ask, not what they ask. Interventions should therefore aim at the decision to ask.

The third contribution is the mechanism. Evaluation apprehension, measured directly in Study 2, carried the effect of the status cue, and anonymity, which removes the basis for evaluation, weakened the effect in Study 1. Both results point to the requester's concern about being judged.

\subsection{Practical implications}

Taken together, these results imply that organizations that display rank in their collaboration tools will receive fewer requests for help from their employees. Channels that detach a question from the person asking it, such as the platform in Study 3, offer one remedy.

\subsection{Limitations}

Studies 1 and 2 used simulated helpers whose replies were prepared in advance, and participants knew the task was a study. Study 3 took place in one firm and one quarter, and the supervisor ratings are the firm's own. Future work could test the effect with real senior colleagues and with outcomes beyond ratings.

\begin{thebibliography}{9}
\setlength{\itemsep}{0pt}

\bibitem[Abadie et al.(2023)]{abadie2023}
Abadie, A., Athey, S., Imbens, G. W., and Wooldridge, J. M. (2023). When should you adjust standard errors for clustering? \textit{Quarterly Journal of Economics}, 138(1), 1--35.

\bibitem[Bruhn and McKenzie(2009)]{bruhn2009}
Bruhn, M., and McKenzie, D. (2009). In pursuit of balance: Randomization in practice in development field experiments. \textit{American Economic Journal: Applied Economics}, 1(4), 200--232.

\bibitem[Lakens(2022)]{lakens2022}
Lakens, D. (2022). Sample size justification. \textit{Collabra: Psychology}, 8(1), 33267.

\bibitem[Nosek et al.(2018)]{nosek2018}
Nosek, B. A., Ebersole, C. R., DeHaven, A. C., and Mellor, D. T. (2018). The preregistration revolution. \textit{Proceedings of the National Academy of Sciences}, 115(11), 2600--2606.

\end{thebibliography}

\clearpage
\appendix
\section*{Appendix A: Materials}

\subsection*{Study 1 helper profiles}

\textit{Senior director condition.} ``Dana Whitlock is a senior director in your division. She has worked in the division's planning group for six years and reports to the chief operating officer. She is known for sharp, demanding feedback. During the task, Dana is available to answer questions.''

\textit{Peer condition.} ``Dana Whitlock is a planning analyst in your division, at the same level as you. She has worked in the planning group for six years. During the task, Dana is available to answer questions.''

\textit{Identified requests.} ``If you ask Dana for a hint, your first name will appear on the request.''

\textit{Anonymous requests.} ``If you ask Dana for a hint, she will see only the text of the request, not your name.''

\textit{Attention item (after the profile).} ``To show that you are reading carefully, please select `Somewhat agree' for this statement.''

\subsection*{Study 2 helper profiles}

\textit{Vice-president condition.} ``Dana Whitlock is the vice president of operations. She has worked in the operations group for six years and reports to the chief executive. During the task, Dana is available by chat to answer questions.''

\textit{Peer condition.} ``Dana Whitlock is an operations analyst, at the same level as you. She has worked in the operations group for six years. During the task, Dana is available by chat to answer questions.''

\subsection*{Study 2 measures}

\textit{Prior help-seeking habit} (1 = strongly disagree, 7 = strongly agree), measured before the profile: ``When I am stuck at work or school, I ask someone quickly''; ``I would rather ask than struggle alone''; ``I often ask people I do not know well for help''; ``Asking questions comes easily to me''; ``I ask for help before a problem gets large.''

\textit{Help-seeking intention} (1 = very unlikely, 7 = very likely): ``How likely are you to contact Dana if you get stuck?''; ``How likely are you to ask Dana to check your work?''; ``How likely are you to ask Dana a question you think may be basic?''; ``How likely are you to contact Dana more than once?''

\textit{Drafted message.} ``Please write the message you would send to Dana if you decided to ask her for help with the task.'' Coders rated each message from 1 (vague, no detail about the problem) to 5 (states the problem, what was tried, and what is needed).

\textit{Comprehension question (after the task).} ``According to the task instructions, which column had to match the monthly totals?''

\textit{Evaluation apprehension} (1 = strongly disagree, 7 = strongly agree), measured after the comprehension question: ``I worried that Dana would think less of me if I asked for help''; ``I was concerned about how my questions would make me look''; ``I felt Dana would judge my ability by my questions''; ``Asking Dana for help felt risky.''

\textit{Manipulation check.} ``How senior is Dana relative to you?'' (1 = much more junior, 7 = much more senior).

\clearpage
\section*{Appendix B: Study 2 Registration}

The text below reproduces the registration of Study 2, filed on February 3, 2025, before data collection began, at \url{https://registry.example.org/r/7Q4K2}.

\paragraph{1. Hypotheses.} (a) Participants who read that the available helper is a vice president will report lower help-seeking intention than participants who read that the helper is a peer. (b) Participants in the vice-president condition will send fewer help requests during the task than participants in the peer condition. This prediction is directional. (c) Evaluation apprehension will mediate the effect of condition on help-seeking intention.

\paragraph{2. Design.} Two conditions, between subjects: vice president or peer. Participants will be randomly assigned to one of the two conditions.

\paragraph{3. Sample.} We will recruit 240 participants, 120 per condition, from the Eastmoor University Behavioral Lab participant pool.

\paragraph{4. Measures and order.} Age; prior help-seeking habit (five items); an attention item; helper profile; help-seeking intention (four items; primary outcome); a drafted help request, rated for quality by two coders who do not know the condition (secondary outcome); the spreadsheet task, during which we count help requests sent to the helper (secondary outcome); a comprehension question about the task rules; evaluation apprehension (four items), at the end of the session; perceived seniority of the helper.

\paragraph{5. Exclusions.} We will exclude participants who fail the attention item shown before the manipulation. There will be no other exclusions.

\paragraph{6. Analyses.} Hypothesis (a): analysis of covariance of help-seeking intention on condition, with age and prior help-seeking habit as covariates, two-tailed, alpha = .05. Hypothesis (b): regression of the number of help requests on condition and the same covariates; because the prediction is directional, we will use a one-tailed test, alpha = .05. Hypothesis (c): the indirect effect of condition on help-seeking intention through evaluation apprehension, with the same covariates and a 95 percent percentile bootstrap interval from 5,000 resamples. The quality of the drafted request will be analyzed by analysis of covariance with the same covariates, two-tailed. Help-seeking intention is the only primary outcome; the two secondary outcomes will be reported as secondary and will not be corrected for multiple comparisons.

\end{document}
